% auto-generated by make_tables_v2.py — DO NOT EDIT BY HAND
\begin{tabular}{lccccl}
\toprule
\textbf{流行波次 / 决策场景} & $\tau=0.20$ (刚性生命线) & $\tau=0.35$ (资源调配线) & $\tau=0.50$ (群体干预线) & $\tau=0.70$ (战略规划线) & \textbf{推荐业务预测与公共卫生风控策略} \\
\midrule
COVID-19 Delta 暴发期 & 9.6 周 (14.3 代) & 15.9 周 (23.7 代) & 21.2 周 (31.6 代) & 27.1 周 (40.3 代) & 低离散度波次，中长程外推相对平稳；采用集合预测中位数推进梯次床位筹备。 \\
COVID-19 Omicron 达峰期 & 3.6 周 (8.3 代) & 5.9 周 (13.8 代) & 7.9 周 (18.5 代) & 10.1 周 (23.5 代) & 高传播株，视界收缩；严控长程外推，实行见转折即报警的敏捷应急机制。 \\
COVID-19 JN.1 流行期 & 4.9 周 (9.7 代) & 8.1 周 (16.2 代) & 10.8 周 (21.6 代) & 13.7 周 (27.5 代) & 平缓爬坡期；聚焦 2--4 周内抗病毒药物分发与高危人群重点防护。 \\
流感 2022--23 暴发早期 & 4.8 周 (10.6 代) & 8.1 周 (17.7 代) & 10.8 周 (23.7 代) & 13.8 周 (30.2 代) & 快速早期爬坡季；第 2--3 周完成儿科药物与门急诊运力跨区协同调配。 \\
流感 2024--25 流行季 & 7.0 周 (15.4 代) & 11.7 周 (25.6 代) & 15.7 周 (34.2 代) & 19.9 周 (43.5 代) & 典型季节性流感波次；匹配 1--2 个月疫苗加强接种与聚集管控预警。 \\
RSV 2024--25 流行季 & 18.6 周 (15.5 代)$^\dagger$ & 30.9 周 (25.8 代)$^\dagger$ & 41.5 周 (34.6 代)$^\dagger$ & 52.7 周 (43.9 代)$^\dagger$ & 数学外推根超单流行季自然跨度（16--20 周截断）；转向宏观峰值高度与总负荷评估。 \\
RSV 2025--26 流行季 & 30.3 周 (25.2 代)$^\dagger$ & 50.9 周 (42.4 代)$^\dagger$ & 67.8 周 (56.5 代)$^\dagger$ & 87.0 周 (72.5 代)$^\dagger$ & 极低初始增长波次；放弃长程点位预报，采用流行季宏观总发病包络区间。 \\
\midrule
\textbf{非 RSV 视界总体区间} & \textbf{3.6--9.6 周} & \textbf{5.9--15.9 周} & \textbf{7.9--21.2 周} & \textbf{10.1--27.1 周} & \textbf{公卫应对全景}：由战术刚性应急向战略资源储备梯次推进。 \\
\bottomrule
\end{tabular}
